\section{Mathematical derivations}
\label{sec:proofs}

\subsection{First-visit law and exponent clipping}

Fix a strategy field and the resource field used for imitation. A cooperator copying a defector contributes $+1$ to the defector increment $Z$; the reverse copying event contributes $-1$. All other visits contribute zero. Uniform selection of the first site and its neighbor gives
\begin{equation}
\label{eq:probabilities}
p_+=\frac{1}{4N}\sum_{e\in\mathcal I}F(x_e),
\qquad
p_-=\frac{1}{4N}\sum_{e\in\mathcal I}F(-x_e).
\end{equation}
An undirected interface edge can be sampled in either orientation, each with probability $1/(4N)$ during the first visit. All other visits leave the defector count unchanged. The remaining probability is $1-p_+-p_-$; it is nonnegative because the square lattice has at most $2N$ undirected edges and $p_++p_-=P/(4N)\le1/2$. The relations below follow from the symmetry of the smooth Fermi rule and concern the single visit conditioned on the initial configuration.

To prove complementarity, write $F(x)=e^x/(1+e^x)$ and $F(-x)=1/(1+e^x)$. Their sum is one. Expressing the hyperbolic functions through exponentials also gives $2F(x)-1=\tanh(x/2)$. Substituting into Eq.~\eqref{eq:probabilities} yields $p_+-p_-=J/N$ and $p_++p_-=P/(4N)$. Since $Z\in\{-1,0,1\}$, its mean is $p_+-p_-$ and its second moment is $p_++p_-$. Subtracting the squared mean proves the variance in Eq.~\eqref{main-eq:moments}.

Moreover, $0<F(x)<1$ implies $|2F(x)-1|\le1$. The finite-sum triangle inequality gives
\begin{equation}
\label{eq:drift_bound}
|J|\le\frac14\sum_{e\in\mathcal I}|2F(x_e)-1|
\le\frac P4.
\end{equation}
When $P=0$, all three first-visit moments vanish. These statements require no independence between successive visits: they describe one visit conditioned on its initial state.

For the numerical clipping convention, define $c_M(x)=\max\{-M,\min(M,x)\}$ with $M=30$. It obeys $c_M(-x)=-c_M(x)$, including at the clipping endpoints. Thus $F(c_M(x))+F(c_M(-x))=1$, and the same second-moment identity holds exactly. The exact clipped mean uses
\begin{equation}
\label{eq:clipped}
J^{(M)}=\frac14\sum_{e\in\mathcal I}
\tanh[c_M(x_e)/2]
\end{equation}
in place of the smooth $J$. The identities for the clipped mean use $J^{(M)}$, whereas the smooth diagnostic $J$ is used for the sensitivity derivatives and reported interface-bias plots. These statements are algebraic properties of the clipping convention; they do not equate exact real arithmetic with floating-point simulation behavior.

\subsection{Initial interface geometry}

Immediately after seeding, every cooperator has resource one and every defector has resource zero. The local reaction leaves these values unchanged. If a boundary cooperator has $a$ cooperative neighbors, its resource after diffusion is
\begin{equation}
\label{eq:initial_resource}
r_C^+=1-\frac D4(4-a).
\end{equation}
Consider a neighboring defector with $c$ cooperative neighbors. The initial payoff difference on this edge is therefore
\begin{equation}
\label{eq:initial_gap}
\Pi_D-\Pi_C
=bc-a\!\left[1+\lambda\sigma
 \left(1-\frac D2(4-a)\right)\right]
\end{equation}
on the positive reward branch. Let $H_{ac}$ count interface edges having these two neighbor counts, where $0\le a\le3$ and $1\le c\le4$. Grouping the interface sum gives
\begin{equation}
\label{eq:motifs}
J_{\rm pre}(0)=\frac14
\sum_{a=0}^{3}\sum_{c=1}^{4}
H_{ac}\tanh\!\left[\frac{\Pi_D(c)-\Pi_C(a)}{2K}\right].
\end{equation}
For an isolated defector, $H_{3,4}=4$ and all other counts vanish. Substitution proves Eq.~\eqref{main-eq:isolated}. Its zero-bias temptation is
\begin{equation}
\label{eq:isolated_zero}
b_0(D)=\frac34[1+\lambda\sigma(1-D/2)].
\end{equation}
At the reference coupling, this is $1.125-0.1875D$. It is an initial local payoff balance, not a predicted depletion threshold. Multiple seeds can alter one another's neighbor counts and resource environments, which is why the measured $H_{ac}$ rather than $m$ controls their initial interface bias.

\subsection{Total resource balance and diffusion modes}

Define $u_i$ by Eq.~\eqref{main-eq:resource}. Every site appears four times in the nearest-neighbor sum on the periodic lattice. Hence
\begin{equation}
\label{eq:resource_total}
\sum_i r_i^+
=(1-D)\sum_i u_i+\frac D4\sum_i\sum_{j\sim i}u_j
=\sum_i u_i.
\end{equation}
Expanding the local increment yields
\begin{equation}
\label{eq:local_balance}
u_i-r_i=g s_i-\delta r_i+(\delta-g)s_i r_i.
\end{equation}
Summation and division by $N>0$ give the first line of Eq.~\eqref{main-eq:balance}. Substituting $\langle sr\rangle=\rho\bar r+C_{sr}$ and collecting terms proves the covariance form. This derivation retains the spatial correlation exactly. Dropping $C_{sr}$ would be an additional closure assumption.

For a fixed set of sites $\Omega$, the same calculation separates local reaction from transport:
\begin{equation}
\label{eq:flux}
\sum_{i\in\Omega}(r_i^+-u_i)
=\frac D4
\sum_{\substack{i\in\Omega,\ j\notin\Omega\\j\sim i}}
(u_j-u_i).
\end{equation}
To obtain this expression, write each diffusion increment as $(D/4)\sum_{j\sim i}(u_j-u_i)$. For an internal pair $i,j\in\Omega$, the two terms are opposite and cancel. Only boundary-crossing pairs remain. The formula concerns a fixed region during the resource substep; a moving strategy domain has additional membership changes. We do not infer a measured front velocity or a flux-controlled growth law from this identity.

For pure diffusion, a lattice Fourier mode with wave vector $(k_x,k_y)$ has multiplier
\begin{equation}
\label{eq:diffusion_mode}
\Lambda_D(k_x,k_y)=
1-D+\frac D2(\cos k_x+\cos k_y),
\end{equation}
with $k_\alpha=2\pi n_\alpha/L$. Each pair of opposite translations contributes $2\cos k_\alpha$ times the mode, giving Eq.~\eqref{eq:diffusion_mode}. The same relation follows for separable real sine and cosine modes by the angle-addition identities. Since each cosine lies in $[-1,1]$,
\begin{equation}
\label{eq:mode_bound}
1-2D\le\Lambda_D\le1,
\qquad |\Lambda_D|\le1.
\end{equation}
For even $L$, the checkerboard mode has multiplier $1-2D$, which is negative when $D>1/2$. Diffusion can therefore reverse a mode's sign without increasing its magnitude. For perturbations about a homogeneous resource equilibrium, local reaction multiplies this factor by $1-g$ in an all-cooperator field or by $1-\delta$ in an all-defector field. For a nonuniform strategy pattern, multiplication by $A_s$ mixes the Fourier modes. A single diffusion eigenvalue does not determine the duration of a resource intervention in the coupled system.

\subsection{Existence, uniqueness, and accuracy of the adapted field}

For fixed $s$, the local maps are $f_C(r)=(1-g)r+g$ and $f_D(r)=(1-\delta)r$. If $r\in[0,1]$ and $0\le g,\delta\le1$, both maps remain in $[0,1]$. Diffusion is a convex average, so $T_s$ maps the closed cube $[0,1]^N$ into itself.

For two fields, the local differences satisfy $|f_{s_i}(r_i)-f_{s_i}(\widetilde r_i)|\le q|r_i-\widetilde r_i|$, where $q=\max\{1-g,1-\delta\}$. Nonnegative averaging weights with unit row sum give
\begin{equation}
\label{eq:average}
\begin{aligned}
|T_s(r)_i-T_s(\widetilde r)_i|
&\le\sum_j (M_D)_{ij}q|r_j-\widetilde r_j|\\
&\le q\|r-\widetilde r\|_\infty.
\end{aligned}
\end{equation}
Taking the maximum and applying induction proves Eq.~\eqref{main-eq:contraction}.

If $0<g,\delta\le1$, then $q<1$. Starting from any $r^{(0)}$ in the cube, set $r^{(n+1)}=T_s(r^{(n)})$. Consecutive differences obey
\begin{equation}
\label{eq:iterate_gap}
\|r^{(n+1)}-r^{(n)}\|_\infty
\le q^n\|r^{(1)}-r^{(0)}\|_\infty.
\end{equation}
For $m>n$, sum these differences to obtain
\begin{equation}
\label{eq:cauchy}
\|r^{(m)}-r^{(n)}\|_\infty
\le\frac{q^n}{1-q}\|r^{(1)}-r^{(0)}\|_\infty.
\end{equation}
The finite-dimensional real space is complete, so the sequence converges. The cube is closed, hence its limit belongs to the cube. Continuity of $T_s$ makes the limit a fixed point. Two fixed points at distance $d$ satisfy $d\le qd$, forcing $d=0$. This proves existence, uniqueness, and convergence to the physical adapted field through the Banach contraction theorem.

The affine form also displays the equilibrium equation explicitly:
\begin{equation}
\label{eq:equilibrium_linear}
(I-M_DA_s)r_{\rm eq}=gM_Ds.
\end{equation}
Since $\|M_DA_s\|_\infty\le q<1$, its inverse can be represented by the convergent geometric operator series. This is an alternative expression for the same fixed point, not a separate physical approximation.

For any trial field $r$, the triangle inequality and contraction give
\begin{equation}
\label{eq:residual_proof}
\begin{aligned}
\|r-r_{\rm eq}\|_\infty
&\le\|r-T_s(r)\|_\infty
  +\|T_s(r)-T_s(r_{\rm eq})\|_\infty\\
&\le\|r-T_s(r)\|_\infty+q\|r-r_{\rm eq}\|_\infty.
\end{aligned}
\end{equation}
Rearranging proves Eq.~\eqref{main-eq:equilibrium_error}. If an iteration stops when $\|r^{(n+1)}-r^{(n)}\|_\infty\le\varepsilon$ and returns the new iterate, one further use of contraction gives
\begin{equation}
\label{eq:returned_error}
\|r^{(n+1)}-r_{\rm eq}\|_\infty
\le\frac{q}{1-q}\varepsilon.
\end{equation}
For $q=0.9$ and a stopping tolerance of $10^{-13}$, the real-arithmetic returned-iterate bound is $9\times10^{-13}$, provided termination occurs by the tolerance criterion. The implementation also imposes a 500-step cap; reaching that cap alone does not certify the same accuracy. These bounds do not account for floating-point roundoff.

The diagnostics $\eta$ and $J_{\rm eq}$ in Eq.~\eqref{main-eq:adaptation} use this fixed-strategy field. In the coupled system, both $s$ and its equilibrium field change. Their definitions do not imply that $\eta$ or $J-J_{\rm eq}$ is a sufficient coordinate for future depletion.

\subsection{Sensitivity through the resource substep}

For one interface edge incident on a boundary cooperator $u$, write
\begin{equation}
\label{eq:edge_derivative}
J_e=\frac14\tanh z_{uv},
\qquad
\frac{\partial z_{uv}}{\partial r_u}
=-\frac{\lambda\sigma k_C(u)}{K}.
\end{equation}
On the positive reward branch, $R'=2\lambda\sigma$. Applying the derivative $(\tanh z)'=1-\tanh^2z$ and summing all incident defector edges proves Eq.~\eqref{main-eq:sensitivity}. The strategy pattern and neighbor counts are held fixed. A defector's own resource does not enter its payoff, so its direct derivative in $J(s,r)$ vanishes.

The intervention is applied to $r$, whereas imitation uses $T_s(r)$. For a direction $h$, the affine resource map gives the exact variation
\begin{equation}
\label{eq:resource_variation}
T_s(r+\epsilon h)-T_s(r)=\epsilon M_DA_sh.
\end{equation}
At a smooth reference point, the chain rule yields
\begin{equation}
\label{eq:chain}
\begin{aligned}
\left.\frac{d}{d\epsilon}J(s,T_s(r+\epsilon h))\right|_0
&=-w^\mathsf{T}M_DA_sh\\
&=-\sum_u a_u(M_D^\mathsf{T}w)_u h_u.
\end{aligned}
\end{equation}
The second line follows by exchanging the finite sums; it does not require independent resources. Symmetry of periodic nearest-neighbor diffusion gives $M_D^\mathsf{T}=M_D$, producing $W_u$ in Eq.~\eqref{main-eq:effective_weight}. The Taylor remainder is quadratic in the maximum norm, for fixed $N$, strategies, and model parameters with $K>0$, along perturbations whose entire segment remains on the positive reward branch. Indeed, $|\tanh''z|\le2$ bounds the remainder of an edge contribution by $(\Delta z_{uv})^2/4$. Convex diffusion and the local slopes give $|\Delta z_{uv}|\le |\lambda\sigma|k_C(u)q\|h\|_\infty/K$. Summing over the fixed interface proves the stated order, with a constant no larger than $9P(\lambda\sigma)^2q^2/(4K^2)$ because a boundary cooperator has at most three cooperative neighbors. This constant is not uniform as $K\to0$ or the interface grows. Finite permutations are evaluated with the nonlinear expression.

This derivation also explains why moving resource among defectors can affect the next strategy-update cycle. Their direct weights $w_u$ vanish, but a neighboring boundary cooperator can contribute to $W_u$ through diffusion. There is no general geometry-independent ratio between the two intervention responses.

\subsection{Rearrangement inequality and invariants}

Let $\mathcal B$ be the fixed set of boundary cooperators. On this set, $a_u=1-g$ is common, so maximizing the linear protective response amounts to maximizing $\sum_{u\in\mathcal B}W_u r_{\pi(u)}$ over permutations $\pi$. Order the weights as $W_1\le\cdots\le W_m$ and the resource values as $r_1\le\cdots\le r_m$. The rearrangement inequality states
\begin{equation}
\label{eq:rearrangement}
\sum_{i=1}^m W_i r_{m+1-i}
\le \sum_{i=1}^m W_i r_{\pi(i)}
\le \sum_{i=1}^m W_i r_i.
\end{equation}
For completeness, consider two weights $W_i\le W_j$ and values $r_i\le r_j$. The gain from aligning their orders is
\begin{equation}
\label{eq:exchange}
(W_i r_i+W_j r_j)-(W_i r_j+W_j r_i)
=(W_j-W_i)(r_j-r_i)\ge0.
\end{equation}
Whenever an assignment is not ordered, exchanging an inverted pair cannot decrease the objective. Repeated exchanges reach the same-order assignment in finitely many steps, proving the upper bound. Reversing the ordering proves the lower bound. Ties yield zero exchange gain and do not change the result.

For any function $\phi$ of a resource value, bijectivity of $\pi$ gives
\begin{equation}
\label{eq:permutation}
\sum_{u\in\mathcal B}\phi(r_{\pi(u)})
=\sum_{u\in\mathcal B}\phi(r_u).
\end{equation}
Taking $\phi(r)=r$ proves preservation of the boundary total and mean. Taking indicator functions proves preservation of the empirical resource distribution. Sites outside $\mathcal B$ are unchanged, so the global distribution and total are also preserved. The nonnegative factor $1-g$ retains the ordering of the linear objective; at $g=1$ it vanishes. The full nonlinear $J_{\rm pre}$ and the continuation fate need not share these extrema.

\subsection{Translation preserves resource autocorrelation}

On the finite periodic lattice, the map $i\mapsto i-\mathbf d$ is a bijection and commutes with every other translation. Since it also preserves the spatial mean, the autocorrelation of a translated field is
\begin{equation}
\label{eq:translation_autocorrelation_proof}
\begin{aligned}
\Gamma_{\mathcal T_{\mathbf d}r}(\mathbf a)
&=\frac1N\sum_i(r_{i-\mathbf d}-\bar r)
                  (r_{i+\mathbf a-\mathbf d}-\bar r)\\
&=\frac1N\sum_j(r_j-\bar r)(r_{j+\mathbf a}-\bar r)
=\Gamma_r(\mathbf a).
\end{aligned}
\end{equation}
The substitution $j=i-\mathbf d$ proves the result at every lag. The covariance with an untranslated strategy field has no analogous invariance, because the strategy factor changes under the same reindexing.

\subsection{Cumulative response and finite event counts}

The cumulative relation in Eq.~\eqref{main-eq:cumulative_response} follows by summing consecutive differences of $\mathbb E_a[n_D(u)]-\mathbb E_0[n_D(u)]$. Interior terms cancel and the initial term vanishes because the strategies initially coincide. Since $0\le n_D\le N$ and the sum is finite, expectation and summation commute without an independence assumption.

For the empirical reach decomposition, let $\mathcal B_T$ and $\mathcal R_T$ contain the realizations that collapse and reach the selected population threshold by $T$. Here $\mathcal B_T\subseteq\mathcal R_T$. Multiplying the reaching fraction by the conditional collapse fraction cancels $|\mathcal R_T|$, provided this count is positive, proving the sample counterpart of Eq.~\eqref{main-eq:reach_decomposition}. The two reported reaching counts are nonzero. Similarly, the eventual collapsed set contains $\mathcal B_T$ and is contained in $\mathcal B_T\cup\mathcal U_T$. Monotonicity of finite cardinality gives Eq.~\eqref{eq:unresolved_bound}.

Subtracting the two resource balances with equal initial $\rho$ and $\bar r$ proves Eq.~\eqref{main-eq:translation_balance}. Solving the same balance for $C_{sr}$ when $\delta\ne g$ proves Eq.~\eqref{eq:reconstruct_covariance}.
